\fichatitulo{19 · INSTITUCIONAL}{Guia · julho de 2024}{Diretrizes de IA em parlamentos}
{\small Fitsilis; von Lucke; De Vrieze (eds.). \textit{Guidelines for AI in Parliaments}. Guia · julho de 2024. \href{https://www.agora-parl.org/sites/default/files/agora-documents/English%20-%20WFD%20AI%20Guidelines%20for%20Parliaments.pdf}{Fonte consultada}.\par}

\campo{Problema e objetivo}
Que critérios devem orientar a adoção parlamentar de IA?

\campo{Principal contribuição · síntese interpretativa}
Organiza 40 diretrizes em seis áreas, relacionando governança, direitos, desenho de sistemas e capacitação.

\campo{Método / natureza da evidência}
Documento orientador construído com contribuições de especialistas; o recorte focal trata de impacto, confiabilidade e aceitação.

\campo{Resultado ou argumento principal}
Recomenda testes, monitoramento, correção de falhas e acordo com interessados sobre a precisão mínima antes da aceitação do sistema.

\campo{Excertos-chave · seleção literal}
\begin{itemize}
\excerto{1}{is dependent on the specific application and intended use}{Diretriz 5.7.}
\excerto{2}{There are no known examples.}{Diretriz 5.3.}
\end{itemize}

\campo{Limitações e análise crítica}
O próprio guia informa ausência de exemplos conhecidos para algumas recomendações. Análise da ficha: diretriz não demonstra eficácia empírica nem equivale a certificação ou obrigação legal.

\campo{Aproveitamento na dissertação · proposta}
Metodologia: discutir critérios de aceitação por tarefa e custos de falsos positivos e negativos.

\registro{Fonte consultada nesta preparação: Apresentação e diretrizes 4.4, 5.3 e 5.7. A chave bibliográfica registra 2023; o PDF consultado é de 2024. Chave: \texttt{wfd2023guidelinesAiParliaments}.}
